% gauge_sector.tex
% The gauge-sector channel: the induced Maxwell action's two anisotropic blocks
% (permittivity eps=P={1,4,4}, inverse permeability mu^-1=Q_B={4,4,16}), the
% adjugate relation, the null birefringence verdict, and the (separate) common-mode
% vacuum-speed anisotropy.
%
% STATUS (audit_table.tex is authority):
%   birefringence verdict (cancels) -> PASS  (null split; engine blocks + solver)
%   common-mode vacuum-speed anisotropy -> PART  (single-domain factor ~2 vs O_h)
% Framing: geometric STRUCTURE (from-engine holonomy) x scalar MAGNITUDE (deferred
% 1/g^2). The single-object dynamical (Tr ln T) tensor extraction is a STATED OPEN
% item (see the remark at the end). Derivation + adjugate theorem: Paper~VIII
% \cite{menendez2026electricblock}. Engine extraction + verdict: exp_03_geometric_blocks,
% exp_03_dispersion_core; notes/exp_03_dynamical_confirmation_plan.md.

\label{sec:gauge_sector}

The same bipartite geometry that fixes the kinematic optical axis also fixes a
\emph{second} anisotropy, in the gauge sector.  Integrating the matter tokens out
of the A=1 dynamics in a background electromagnetic field induces an effective
Maxwell action whose energy density is the anisotropic quadratic form
\begin{equation}
    u = \tfrac12\bigl(\mathbf{E}^{T}\boldsymbol{\varepsilon}\,\mathbf{E}
        + \mathbf{B}^{T}\boldsymbol{\mu}^{-1}\mathbf{B}\bigr),
    \label{eq:induced_action}
\end{equation}
with a permittivity tensor $\boldsymbol{\varepsilon}$ and an inverse-permeability
tensor $\boldsymbol{\mu}^{-1}$ that are \emph{not} isotropic.  Both are uniaxial
about the same optical axis $\hat{\mathbf{n}}_* = (1,1,-1)/\sqrt3$ that governs the
kinematic channel (Sec.~\ref{sec:kinematic_channel}).  We report the two tensors as
they are read directly off the engine's hop geometry, establish that the induced
birefringence \emph{cancels}, and separate that (settled) null result from the
distinct---and open---question of the common-mode vacuum speed.  The overall scale
of Eq.~\eqref{eq:induced_action} carries the one undetermined coupling $1/g^2$ that
Paper~I leaves open; throughout, we display anisotropic \emph{structure} and treat
that isotropic \emph{magnitude} as deferred.  The analytic derivation of the blocks
and of the cancellation theorem is the companion Paper~VIII
\cite{menendez2026electricblock}; the present section is the independent from-engine
confirmation.

\subsection{The two blocks from the hop geometry}
\label{subsec:blocks}

Write the three RGB hop vectors as $\mathbf{V}_1=(1,1,1)$, $\mathbf{V}_2=(1,-1,-1)$,
$\mathbf{V}_3=(-1,1,-1)$ (as in Sec.~\ref{sec:kinematic_channel}).  The two blocks
are the two natural quadratic forms these vectors build:
\begin{equation}
    \boldsymbol{\varepsilon} = P = \sum_{a} \mathbf{V}_a\mathbf{V}_a^{T},
    \qquad
    \boldsymbol{\mu}^{-1} = Q_B
        = \sum_{a<b} (\mathbf{V}_a\times\mathbf{V}_b)(\mathbf{V}_a\times\mathbf{V}_b)^{T}.
    \label{eq:blocks_def}
\end{equation}
The electric block sums the hop vectors themselves; the magnetic block sums their
oriented \emph{plaquette} areas $\mathbf{V}_a\times\mathbf{V}_b$.  Explicitly,
\begin{equation}
    \boldsymbol{\varepsilon} = \begin{pmatrix} 3&-1&1\\-1&3&1\\1&1&3 \end{pmatrix},
    \quad \text{eig}\{1,4,4\};
    \qquad
    \boldsymbol{\mu}^{-1} = \begin{pmatrix} 8&4&-4\\4&8&-4\\-4&-4&8 \end{pmatrix},
    \quad \text{eig}\{4,4,16\}.
    \label{eq:blocks_explicit}
\end{equation}
Both share the eigenvector $\hat{\mathbf{n}}_*=(1,1,-1)/\sqrt3$, but with opposite
sense: the optical axis is the \emph{suppressed} ($\varepsilon=1$) direction of the
permittivity and the \emph{enhanced} ($\mu^{-1}=16$) direction of the inverse
permeability.  The origin of the axis is geometric: $\hat{\mathbf{n}}_*$ is the one
cube body-diagonal \emph{absent} from the hop set $\{\mathbf{V}_a\}$, so the hop set
breaks the cubic group $O_h$ to the trigonal $D_{3d}$ about $\hat{\mathbf{n}}_*$
\cite{menendez2026electricblock}.

Each block is read off the running engine, not asserted.  The magnetic block is the
spatial plaquette holonomy of a uniform field: with the Peierls link phase
$\exp(i\,\tfrac12(\mathbf{A}(\mathbf{x})+\mathbf{A}(\mathbf{x}-\mathbf{V}))\cdot\mathbf{V})$
threaded through the hop, the smallest closed loop spanned by $\mathbf{V}_a,\mathbf{V}_b$
carries flux $(\mathbf{V}_a\times\mathbf{V}_b)\cdot\mathbf{B}$, and summing its square
over the RGB planes reproduces $Q_B$ (the dcl-core \texttt{exp\_04} route).  The
electric block has no purely spatial holonomy---the field enters the tick rule
\emph{on-site}, through $\delta\phi(\mathbf{x}) = \omega + V(\mathbf{x})$ with
$V=-\mathbf{E}\cdot\mathbf{x}$---so it is read from the \emph{temporal} plaquette:
seeding a uniform real state and taking one real step recovers
$\delta\phi = \omega + V$ from $\mathrm{Im}\,\psi_R^{\text{new}}/\psi_R
= \sin(\delta\phi/2)$, and the loop
$\Theta_a(\mathbf{x}) = \delta\phi(\mathbf{x}) - \delta\phi(\mathbf{x}+\mathbf{V}_a)
= \mathbf{V}_a\cdot\mathbf{E}$ (the mass term $\omega$ cancels around the loop)
gives $\sum_a\Theta_a^2 = \mathbf{E}^{T}P\,\mathbf{E}$
\cite{menendez2026electricblock}.  The experiment
\texttt{exp\_03\_geometric\_blocks} performs both extractions from one engine and
recovers Eq.~\eqref{eq:blocks_explicit} with the optical axis in the expected
(suppressed/enhanced) sense; the on-site coupling reproduces $\omega+V$ to
$4\times10^{-19}$ (unit tick weight), and the extracted $P$ is independent of
$\omega$ to $7\times10^{-15}$ (the mass cancels), so the number genuinely exercises
\texttt{HopOperator.step}.

\subsection{The covariant relation and the null birefringence verdict}
\label{subsec:verdict}

The two blocks are not independent.  For \emph{any} three vectors, the sum of
squared plaquette areas is the adjugate of the sum of squared vectors,
\begin{equation}
    Q_B = \sum_{a<b}(\mathbf{V}_a\times\mathbf{V}_b)(\mathbf{V}_a\times\mathbf{V}_b)^{T}
        = \operatorname{adj}\!\Bigl(\sum_a\mathbf{V}_a\mathbf{V}_a^{T}\Bigr)
        = \operatorname{adj}(P),
    \label{eq:adjugate}
\end{equation}
so $\boldsymbol{\mu}^{-1} = \operatorname{adj}(\boldsymbol{\varepsilon})$ is a
structural identity, not a coincidence of the octahedral lattice
\cite{menendez2026electricblock}.  The two engine-extracted blocks satisfy it: they
commute to $4\times10^{-15}$ and $\operatorname{adj}(\boldsymbol{\varepsilon})$
matches $\boldsymbol{\mu}^{-1}$ to $5\times10^{-15}$
(\texttt{exp\_03\_geometric\_blocks}).  Equation~\eqref{eq:adjugate} is exactly the
condition that removes the birefringence.

\begin{proposition}[Null gauge-sector birefringence]
\label{prop:null_split}
For a medium with $\boldsymbol{\mu}^{-1} = \operatorname{adj}(\boldsymbol{\varepsilon})$,
the photon dispersion factorizes as
\begin{equation}
    \det\!\bigl(K\,\boldsymbol{\mu}^{-1}K + \omega^2\boldsymbol{\varepsilon}\bigr)
    = \det(\boldsymbol{\varepsilon})\,\omega^2\,
      \bigl(\omega^2 - \mathbf{k}^{T}\boldsymbol{\varepsilon}\,\mathbf{k}\bigr)^2 ,
    \qquad K = [\hat{\mathbf{k}}]_\times ,
    \label{eq:dispersion_factored}
\end{equation}
a \emph{double} transverse root for every propagation direction.  Both
polarizations share $\omega^2 = \mathbf{k}^{T}\boldsymbol{\varepsilon}\,\mathbf{k}$,
so the split
$\Delta(\omega^2)\propto |\varepsilon_a\mu^{-1}_a - \varepsilon_p\mu^{-1}_p|
= |\det\boldsymbol{\varepsilon}-\det\boldsymbol{\varepsilon}| = 0$
vanishes.  The cancellation is invariant under independent rescaling of
$\boldsymbol{\varepsilon}$ and $\boldsymbol{\mu}^{-1}$, hence immune to the deferred
$1/g^2$, and holds for any hop vectors \textup{\cite{menendez2026electricblock}}.
\end{proposition}

Feeding the two \emph{engine-extracted} blocks of Eq.~\eqref{eq:blocks_explicit}
into the plane-wave dispersion of Proposition~\ref{prop:null_split} confirms this
numerically: the experiment \texttt{exp\_03\_dispersion\_core} solves the
generalized eigenproblem $-K\boldsymbol{\mu}^{-1}K\,\mathbf{E}
= v^2\boldsymbol{\varepsilon}\,\mathbf{E}$ and finds the two transverse speeds equal
to $\max|v_+-v_-| = 1.6\times10^{-15}$ over $2000$ directions, while a deliberately
detuned control (breaking Eq.~\eqref{eq:adjugate}) returns a split of $0.5$---so the
null is a genuine cancellation, not an insensitive estimator.  This discharges the
gauge-sector \emph{birefringence} audit row to \texttt{PASS}
(Table~\ref{tab:audit}).  Independently, the same operator anisotropy \emph{survives}
the A=1 token renormalization: driving the engine's full multi-tick dynamics with a
background field, the induced second-order response is uniaxial about
$\hat{\mathbf{n}}_*$ with the correct opposite senses ($\mathbf{B}$ axis-enhanced,
$\mathbf{E}$ axis-suppressed), so the cancellation is a property of the renormalized
observables, not only of the tree-level operator.

\subsection{The common-mode vacuum speed}
\label{subsec:common_mode}

The null split does \emph{not} make the gauge sector isotropic.  The shared
dispersion of Proposition~\ref{prop:null_split},
$\omega^2 = \mathbf{k}^{T}\boldsymbol{\varepsilon}\,\mathbf{k}$, is
direction-dependent: the common-mode phase speed
$v^2(\hat{\mathbf{k}}) = \hat{\mathbf{k}}^{T}\boldsymbol{\varepsilon}\,\hat{\mathbf{k}}$
ranges over the eigenvalues of $\boldsymbol{\varepsilon}$, i.e.\ $v^2\in[1,4]$---a
factor-$\sim2$ anisotropy of the vacuum speed about $\hat{\mathbf{n}}_*$, affecting
\emph{both} polarizations equally.  The dispersion solver on the engine blocks
returns $v\in[1.00,2.00]$ for a single lattice.  This is a speed anisotropy, not
birefringence, and its observational status is genuinely open.

\begin{remark}[Single domain versus $O_h$ restoration]
\label{rem:oh_restoration}
The factor-$\sim2$ anisotropy is the response of a \emph{single} trigonal domain.
Averaging the blocks over the four ways of choosing three of the four cube
body-diagonals as hop directions gives
$\langle\boldsymbol{\varepsilon}\rangle = 3\mathbb{I}$ and
$\langle\boldsymbol{\mu}^{-1}\rangle = 8\mathbb{I}$, restoring $O_h$ and removing the
anisotropy entirely (\texttt{exp\_03\_geometric\_blocks}: residual $10^{-16}$).  But
the physical lattice has \emph{one fixed} hop set---its optical axis is a specific
body-diagonal---so this four-domain average is a fictitious ensemble the lattice does
not realize, and there is no dynamical mechanism that performs it.  Whether the
common-mode anisotropy is therefore a real prediction (single domain) or is washed
out (restored) is the isotropy-restoration question this paper leaves at
\texttt{PART}; it is shared with the companion Paper~VIII
\cite{menendez2026electricblock} and is distinct from the null birefringence, which
holds in \emph{every} domain via Eq.~\eqref{eq:adjugate} regardless of the answer.
\end{remark}

\subsection{Scope: structure, magnitude, and the deferred dynamical tensor}
\label{subsec:gauge_scope}

The results above are the anisotropic \emph{structure} of the induced action
(Eq.~\eqref{eq:blocks_explicit}), computed from the engine hop geometry, together
with the exact cancellation it implies.  The isotropic \emph{magnitude}---the single
scalar $1/g^2$ multiplying Eq.~\eqref{eq:induced_action}---is the coupling Paper~I
leaves undetermined and is not a prediction of the vacuum speed here.  Because the
birefringence verdict is prefactor-independent (Proposition~\ref{prop:null_split}),
that omission does not weaken it.

\begin{remark}[The dynamical tensor extraction is a stated open item]
\label{rem:trlnt_open}
The blocks of Eq.~\eqref{eq:blocks_def} are obtained from the hop \emph{geometry}
(the plaquette holonomies), which is the tree-level induced action and is exact at
any lattice size.  One can also ask for the tensors as a genuinely \emph{dynamical},
momentum-space fermion-loop sum---the second-order response of the physical band of
the one-period transfer operator $T(\mathbf{k})$ summed over the Brillouin zone
($\operatorname{Tr}\ln T$).  We do not carry that computation out here, and we do not
claim it.  A naive band sum is not merely hard but Ward-\emph{divergent}: along the
optical axis the physical band is exactly flat (all three hops project equally onto
$\hat{\mathbf{n}}_*$, the $D_{3d}$ symmetry), so the intraband denominator
$\lambda(\mathbf{k})-\lambda(\mathbf{k}\pm\mathbf{q})$ vanishes and the
occupation-cancelled (diamagnetic) seagull must be subtracted for a finite answer;
once regularized, the naive current--current object returns the \emph{wrong} tensor
(uniaxial about $(1,-1,0)$, not the optical axis), and the correct assembly is the
orbital-susceptibility (Berry-curvature plus quantum-metric) combination.  This
divergence and its wrong-axis outcome were diagnosed in Paper~VIII
\cite{menendez2026electricblock} and reproduced here in an exploratory band-sum
probe; the geometric route of Sec.~\ref{subsec:blocks} sidesteps them entirely.
Since the geometric and fermion-loop routes
compute the \emph{same} induced action and must agree---and the geometric route
already yields Eq.~\eqref{eq:blocks_explicit}---the dynamical extraction is a matter
of methodological completeness, not of deciding the verdict, and is deferred to
follow-on work.  As with the kinematic doubler accounting
(Remark~\ref{rem:no_new_modes}), nothing in the birefringence verdict depends on it.
\end{remark}
